\chapter{The data manager}\label{ch:data-manager}

The data manager shows every element of a type in one editable table, the way
PowerFactory's data manager or a spreadsheet does. Use it to check the data of a whole
network at a glance, or to change many elements without clicking each one on the canvas.

\begin{figure}[H]
  \centering\shot[\linewidth]{data-manager.png}
  \caption{The lines of the load flow example in the data manager. The \gui{r0 / r}
    column and the ones after it are greyed out: those fields only apply when \gui{Z0
    data} is \code{factors}.}
\end{figure}

\section{Opening it}

The data manager shares the bottom panel with \gui{Results}: click its tab, or use
\menu{View > Data manager} (\keys{\ctrl + \shift + D}). The menu entry brings the panel
forward on the type of the selected element.

\section{The tables}

There is one tab per component type, in library order, titled with the number of
elements of that type: \gui{Bus (8)}, \gui{Line / Cable (9)}, \dots\ Each row is an
element, in network order. The first column is its id, and the others are its fields,
with the same labels and units as the property editor. Hover a column header to read the
field's description.

\begin{itemize}
  \item \textbf{Greyed-out cells} belong to fields the element's mode does not use, such as
    a load's \gui{cos phi} when it is given by $P$ and $Q$, or a line's zero-sequence
    factors when its \gui{Z0 data} is \code{absolute}. They are still editable.
  \item \textbf{Tinted cells} hold a value that comes from an expression (a binding).
    Hover the cell to see the expression.
  \item \textbf{Bus columns} show the bus's name.
\end{itemize}

\section{Editing}

Double-click a cell, or select it and start typing. Cells take the same input as the
property editor:

\begin{description}
  \item[Numbers and expressions] A number is stored as it is. An expression over the
    project variables becomes a binding. An empty cell removes the binding, and makes
    the value empty for fields that accept that, such as \gui{U nominal} or \gui{Q max}.
  \item[Choices] are picked from a list.
  \item[Check boxes] (\gui{In service}, \gui{Slack}, \dots) toggle with a click.
  \item[Buses] are picked from the list of buses, or typed by name or id. Pick
    \gui{(none)} to disconnect the terminal.
\end{description}

Every edit is one undoable step. The canvas and the property editor follow it at once,
and the results become stale. An edit that is not valid, such as text in a numeric field,
a negative length or a bus that does not exist, is refused. The reason appears in the
status bar.

\section{Sorting and filtering}

\begin{itemize}
  \item Click a column header to sort by that column; click it again to reverse the order.
  \item Type in the \gui{Filter} box to keep the rows where any cell contains the text,
    case aside: \code{L2} keeps L23, L27 and L28.
\end{itemize}

Sorting and filtering only change the view, never the network.

\section{Finding an element on the diagram}

Double-click an element's id to select it on the canvas, centre the diagram on it and
show its properties. \menu{Edit > Find element...} (\keys{\ctrl + F}) does the same from
anywhere (\cref{sec:find}).
